\begin{theorem}[Connected exterior cell contour]%
    \label{thm:peps_integer_boundary_contour_connected}
    \lean{TNLean.PEPS.integerExteriorCollarGraph_connected_of_isSimplyConnected}
    \leanok
    \uses{def:peps_integer_exterior_collar_graph,
        thm:peps_integer_cell_exterior_connected,
        thm:peps_integer_cell_contact_connectivity,
        thm:peps_integer_contour_ray_potential}
    Let $A\subset\mathbb Z^2$ be finite and suppose that the actual union
    $P_A=\bigcup_{a\in A}(a+[-1/2,1/2]^2)$ is simply connected.
    Let $B_A$ consist of the unoccupied centers whose coordinate distances
    from some occupied center are at most one. Join four-neighbor centers
    in $B_A$ whenever one occupied center is within coordinate distance
    one of both. Then this graph is connected.
    This is auxiliary exterior-collar geometry for
    \cite[proof of Theorem~6.9]{Schuch2010PEPS}, local source lines
    1935--2076. Shared corners of occupied cells are allowed; simple
    connectedness refers to the genuine closed-cell union.
\end{theorem}

\begin{proof}\leanok
    Orient each exposed cell edge with its occupied cell on the left.
    At an edge endpoint, continue with a right turn when the diagonal
    cell is occupied, straight ahead when only the forward cell is
    occupied, and a left turn otherwise. The reverse rule makes this
    successor a permutation of the finite exposed-edge set.

    Fix one successor orbit and assign weight one to its edges and zero
    to the remaining edges. The resulting finite flow is conserved.
    Its horizontal-ray winding potential is constant on missing cells,
    since their four-neighbor graph is connected, and vanishes far from
    the occupied cells. It is also constant on occupied cells: the edge
    identities give equality across shared sides, and the right-turn
    rule gives equality across shared corners. The occupied contact
    graph is connected. Across any exposed edge the winding difference
    is its weight. Thus every exposed edge has the same weight, and all
    belong to the chosen successor orbit.

    Consecutive contour edges have exterior centers that are equal or
    joined by one or two four-neighbor steps with a common occupied-cell
    support. Every missing center in $B_A$ can reach one of these
    exterior centers using the same local support. These walks prove
    graph connectivity. Nonemptiness follows by taking an occupied
    center with maximal first coordinate and the missing center to its east.
\end{proof}
